\documentclass[11pt,a4paper]{article}

\usepackage[utf8]{inputenc}
\usepackage[T1]{fontenc}
\usepackage{amsmath,amssymb,amsthm}
\usepackage{graphicx}
\usepackage{booktabs}
\usepackage{hyperref}
\usepackage{xcolor}
\usepackage{tcolorbox}
\usepackage[margin=25mm]{geometry}
\usepackage{xeCJK}
\setCJKmainfont{Hiragino Mincho ProN}
\setCJKsansfont{Hiragino Sans}

\newtheorem{observation}{観察}
\newtheorem{conjecture}{予想}
\newtheorem{theorem}{定理}
\newtheorem{proposition}{命題}

\title{保型形式物理と Wright Brothers プログラム:\\
統合と新しい貢献領域}
\author{Wright Brothers Research Program}
\date{2026年4月}

\begin{document}
\maketitle

\begin{abstract}
本稿は Wright Brothers プロジェクトを既存の保型形式物理
(Heterotic string, Monstrous/Mathieu Moonshine, 位相弦理論,
BPS 状態数え上げ, S 双対等) と体系的に統合する試みである．
第 1 部で既存の保型形式物理の成果を review し，どの分野で
どの modular object が物理的に解釈されているかを整理する．
第 2 部で\textbf{既存の保型形式物理が達成していない 7 つの空白領域}
を identify する: (i) 観測される低エネルギー SM パラメータ
($\alpha^{-1}, m_e/m_\mu$, etc.) の決定, (ii) 宇宙論定数
$\Omega_\Lambda$ の予測, (iii) Wheeler--DeWitt 量子宇宙論との統合,
(iv) 世代階層の arithmetic 起源, (v) Beilinson 予想の物理的実現,
(vi) 低次元 (4D) 観測量と保型形式の直接 bridge,
(vii) 宇宙論定数問題の 120 桁ズレの解消．
第 3 部で Wright Brothers がこれらの空白にどう貢献しうるかの
具体的提案を行う．特に Wheeler--DeWitt 時間的 DN 境界と
Eisenstein 級数の統合 ($\Omega_\Lambda = 2\pi/9$ の導出)，
Ramanujan $\tau$ と generational 異常の対応，そして
Bernoulli ladder と Heterotic $E_8$ の新しい接続を提示する．
Wright Brothers は「counting theory」ではなく「constant
determination theory」として既存の保型形式物理と相補的に
機能できる可能性を示す．
\end{abstract}

\tableofcontents

\newpage
\section{第 1 部: 既存の保型形式物理の概観}

\subsection{保型形式 (modular form) とは}

$SL_2(\mathbb{Z})$ の上半平面への作用に対して特定の変換則を満たす
正則関数を保型形式と呼ぶ．重さ $2k$ のholomorphic modular form
$f(\tau)$ は:
\begin{equation}
f\left(\frac{a\tau + b}{c\tau + d}\right) = (c\tau + d)^{2k} f(\tau),
\quad \begin{pmatrix} a & b \\ c & d\end{pmatrix} \in SL_2(\mathbb{Z})
\end{equation}

Fourier 展開:
\begin{equation}
f(\tau) = \sum_{n=0}^\infty a_n q^n, \quad q = e^{2\pi i \tau}
\end{equation}

重要な例:
\begin{itemize}
\item Eisenstein 級数 $E_{2k}$ (重さ $2k$, $k \geq 2$, 正則)
\item 判別式 $\Delta$ (重さ 12, cusp form)
\item モジュラー $j$ 不変量 (重さ 0)
\end{itemize}

\subsection{物理への登場経路}

保型形式は物理に以下の経路で登場:

\subsubsection{Heterotic string T$^d$ compactification}

Heterotic string を $d$ 次元トーラスにコンパクト化すると，
partition function は modular invariant でなければならず，
Eisenstein 級数と theta 関数が自然に登場する~\cite{GSW1987}．

\begin{itemize}
\item $E_8 \times E_8$ heterotic: 左移動 gauge 自由度が $E_8 \times E_8$
\item $SO(32)$ heterotic: $SO(32)$ の theta
\item $T^d$ 上の 1-loop partition function: $E_4$, $E_6$ が登場
\end{itemize}

Key identity: $E_8$ 格子のテータ関数
\begin{equation}
\theta_{E_8}(\tau) = \sum_{v \in E_8} q^{|v|^2/2} = E_4(\tau)
\end{equation}

これは Wright Brothers が先行 note~\cite{wb-spec-z}で強調した
「E$_8$ root 数 $= 240 =$ Eisenstein coefficient」の\textbf{数学的根拠}．

\subsubsection{Monstrous Moonshine}

Conway--Norton (1979) 予想，Borcherds (1992) 証明~\cite{Borcherds1992}:
Monster 群 $M$ (有限単純群の最大の散発的群, 位数 $\approx 8 \times 10^{53}$)
の既約表現次元と $j$-関数 Fourier 係数が一致:
\begin{equation}
j(\tau) = \frac{1}{q} + 744 + 196884 q + 21493760 q^2 + \ldots
\end{equation}
$196884 = 196883 + 1$ (最小次元 $+$ 自明表現)．

物理的解釈: Frenkel--Lepowsky--Meurman の Monster vertex operator algebra,
Dixon--Ginsparg--Harvey の $Z_2$ orbifold on Leech lattice.
2 次元共形場理論として実現~\cite{FLM1988}．

\subsubsection{Mathieu Moonshine}

Eguchi, Ooguri, Tachikawa (2010)~\cite{EOT2010}: K3 面の elliptic genus
が\textbf{mock modular form}で，Mathieu 群 $M_{24}$ の表現論と関連．
証明は Gannon (2016)~\cite{Gannon2016}．

物理的解釈: K3 sigma 模型，$N = 4$ 位相的弦．

\subsubsection{Gopakumar--Vafa 不変量}

Calabi-Yau 3-fold 上の位相的弦理論の partition function は，
BPS 状態を数える\textbf{Gopakumar--Vafa 不変量}~\cite{GV1998}を
含む．これらは多くの場合 modular form．

\subsubsection{S-duality と $N = 4$ SYM}

$N = 4$ 超対称 Yang--Mills 理論は S-双対性
\begin{equation}
\tau = \frac{\theta}{2\pi} + \frac{i}{g^2}, \quad \tau \to -\frac{1}{\tau}
\end{equation}
を持ち，分配関数は $SL_2(\mathbb{Z})$ 不変．Vafa--Witten 理論の partition
function も modular form~\cite{VafaWitten1994}．

\subsubsection{Black Hole エントロピーと Jacobi 形式}

Dijkgraaf--Moore--Verlinde--Verlinde (1997)~\cite{DMVV1997}: $D = 4$, $N = 4$
ブラックホールの microstate 数は Igusa cusp form から．
Dabholkar et al.\ (2012)~\cite{Dabholkar2012}: 1/4-BPS ブラックホールの
退化度は\textbf{mock Jacobi form}．

\subsubsection{位相弦と数論}

位相的弦理論の全次数 (all-genus) 計算は，多くの場合 quasi-modular
form で表現される．Bershadsky--Cecotti--Ooguri--Vafa (BCOV) の
\textit{holomorphic anomaly equation}~\cite{BCOV1994}は modular
covariance を要求．

\subsection{まとめ: 既存達成}

\begin{center}
\begin{tabular}{lll}
\toprule
分野 & 保型対象 & 主な物理的内容 \\
\midrule
Heterotic string & $E_4, E_6, E_8$ & 10D partition function \\
Monstrous Moonshine & $j(\tau)$ & Monster VOA, CFT \\
Mathieu Moonshine & mock modular & K3 CFT \\
Gopakumar--Vafa & modular forms & CY$_3$ BPS states \\
S-duality (N=4 SYM) & $SL_2(\mathbb{Z})$ & EM duality \\
BH entropy & Igusa forms & Microstate counting \\
BCOV (位相弦) & quasi-modular & Holomorphic anomaly \\
\bottomrule
\end{tabular}
\end{center}

これら全ては\textbf{数学的に rigorous}で，保型形式物理の
flagship 成果．

\section{第 2 部: 保型形式物理が達成していないこと}

上記の成果は profound だが，以下の 7 つの\textbf{空白領域}が存在する．

\subsection{空白 1: 低エネルギー SM パラメータの決定}

既存の保型形式物理は主に高次元・高エネルギー・位相的文脈で機能する．
\textbf{観測される 4 次元 SM パラメータ}を arithmetic から決定する
試みはほぼ存在しない:

\begin{itemize}
\item 微細構造定数 $\alpha^{-1} = 137.036$: 無予測
\item 電子質量 $m_e = 0.511\,$MeV: 無予測
\item 全ての Yukawa 結合: 無予測
\item CKM 混合: 無予測
\item Higgs 自己結合: 無予測
\end{itemize}

Heterotic string の landscape は $10^{500}$ の真空を許し，具体的な
SM パラメータを指定できない~\cite{Susskind2003}．

\subsection{空白 2: 宇宙論定数 $\Omega_\Lambda$ の予測}

String theory は宇宙論定数問題に対して:
\begin{itemize}
\item \textbf{Landscape}: 多くの真空が存在，anthropic 選択~\cite{Weinberg1987}
\item \textbf{de Sitter vs AdS}: de Sitter 真空の construction は困難
(KKLT~\cite{KKLT2003}, swampland conjecture~\cite{Vafa2005})
\item \textbf{小さな正の値の自然な説明}: 存在せず
\end{itemize}

Moonshine や topological string も $\Omega_\Lambda$ を扱わない．

\subsection{空白 3: Wheeler--DeWitt 量子宇宙論との統合}

保型形式物理は mostly 弦理論文脈で，\textbf{量子宇宙論 (Wheeler--DeWitt,
最小超空間)} との直接的 link は薄い．

量子宇宙論研究 (Hartle--Hawking, Vilenkin, Linde 等) は modular form を
使わない．両分野は独立に発展してきた．

\subsection{空白 4: 世代階層の arithmetic 起源}

SM の 3 世代構造は\textbf{最大の謎}の一つ．既存 approach:
\begin{itemize}
\item Family symmetry (Froggatt--Nielsen, $U(1)_{\rm FN}$ 等)
\item Extra dimensions (Kaluza--Klein 世代)
\item String phenomenology (intersecting branes)
\end{itemize}

どれも numerical prediction は与えない．保型形式からの世代階層導出は
ない．

\subsection{空白 5: Beilinson 予想の物理的実現}

数論では Deligne--Beilinson 予想~\cite{Deligne1979,Beilinson1984}が
zeta 特殊値と K 理論を結ぶ．Bloch--Kato 予想~\cite{BlochKato1990}は
これを精密化．これらは $\zeta(s)$ の整数値が motive の cohomology と
深く関連することを予想する．

\textbf{しかしこの予想の物理的実現は未発見}．もし $\zeta(s)$ 特殊値が
物理定数を\textbf{直接}決めるなら，Beilinson 予想は物理的に具体化
される．保型形式物理はこの方向にはほとんど踏み込んでいない．

\subsection{空白 6: 4D 観測量と保型形式の直接 bridge}

既存の保型形式物理は:
\begin{itemize}
\item 10D/11D で機能 (弦理論)
\item トポロジカルに機能 (位相弦，Moonshine)
\item BPS 状態数として (BH, Gopakumar--Vafa)
\end{itemize}

いずれも\textbf{4D 実験室で測定される現実の物理量}とは遠い．
ミューオン $g-2$ のような 4D QED 観測量が保型形式と関連すると
主張する主流理論はない．

\subsection{空白 7: $10^{123}$ ズレの解消}

宇宙論定数問題の「$10^{123}$ 桁ズレ」を\textbf{定量的に}解消する
保型形式ベースの提案は存在しない．既存の string landscape は
multiverse 的アプローチで「可能な値のうち観測値が anthropic」と
言うだけで，specific な値を予測しない．

\section{第 3 部: Wright Brothers の貢献可能性}

各空白に対し，Wright Brothers プログラムがどう貢献しうるかの
具体案．

\subsection{貢献 1: $\alpha^{-1}$ の arithmetic 決定}

先行 note~\cite{wb-alpha}:
\begin{equation}
\alpha^{-1} = \frac{2}{|B_2|} + \frac{4}{|B_4|} + \text{corrections}
\end{equation}

これは保型形式理論で:
\begin{align}
\frac{2}{|B_2|} &= 12 = \text{(負の方向の) Eisenstein $E_2$ coefficient (正規化)}\\
\frac{4}{|B_4|} &= 120 = \text{Eisenstein $E_4$ coefficient (正規化)}/2
\end{align}

(正確には $-4k/B_{2k}$ が $E_{2k}$ の normalized coefficient)．

\textbf{Wright Brothers の主張}: $\alpha^{-1}$ は Eisenstein 級数係数の
\textbf{線形結合}．これは既存の保型形式物理が行わない「4D 観測量の
arithmetic 決定」．

\textbf{導出方向}: QED の loop 展開で Bernoulli 数が residues として
登場する事実~\cite{Kreimer1999}を利用して，$\alpha^{-1}$ の
Eisenstein 表現を rigorous に導く．これは Kreimer の Hopf 代数アプローチと
直接関連する可能性があり，次の research topic．

\subsection{貢献 2: $\Omega_\Lambda = 2\pi/9$ の導出}

先行 note~\cite{wb-wdw}:
\begin{equation}
\Omega_\Lambda = 8\pi B_2^2 = \frac{2\pi}{9}
\end{equation}

\textbf{新しい意味付け}: $B_2^2$ は重さ 4 の modular form 空間
$M_4 = \mathbb{C} \cdot E_4$ での「$E_2^{\rm double}$」的対象の
leading coefficient (Eisenstein $E_2$ は quasi-modular なので
$E_2^2$ を regularize すると $E_4$ と修正項になる)．

より厳密には:
\begin{equation}
E_2^2 - E_4 = -\frac{12 E_2'}{2\pi i}
\end{equation}
(the Ramanujan identity)．$E_2' = dE_2/d\tau$．

この identity は $B_2^2$ 的量の quasi-modular 性を定量化する．
Wright Brothers の $\Omega_\Lambda \propto B_2^2$ 式は
$E_2^2$ の quasi-modular 依存性と関連する可能性がある．

\textbf{これは新しい観察}で，既存の保型形式物理に無い link を示唆．

\subsection{貢献 3: WDW 量子宇宙論 $+$ Eisenstein 級数}

先行 note~\cite{wb-wdw}で提案された WDW 時間的 DN 境界 $+$
CKN holographic bound → $\Omega_\Lambda = 2\pi/9$ 導出を，
modular form framework で reformulate する:

\begin{proposition}[modular reformulation]
WDW 最小超空間の波動関数 $\Psi(a)$ が時間的 DN 境界を満たすとき，
その零点エネルギーは 1D DN cavity Casimir に等しく，
\begin{equation}
E_0^{\rm DN} = \frac{\hbar c \pi}{48 L} = \frac{\hbar c \pi}{4L}|\zeta_{\neg 2}(-1)|
\end{equation}
係数 $|\zeta_{\neg 2}(-1)| = 1/12 = |B_2|/2$ は Eisenstein $E_2$ の
定数項補正 (quasi-modular) に由来する．
\end{proposition}

これは WDW 量子宇宙論と保型形式の\textbf{初めての統合}．両分野は
これまで独立だった．

\subsection{貢献 4: 世代階層と Ramanujan $\tau$}

先行観察~\cite{wb-spec-z}: ミューオン $g-2$ 異常 $\approx 252 = \tau(3)$
(Ramanujan 判別式の第 3 Fourier 係数)．

\begin{conjecture}[世代-モジュラー対応]
SM の世代 $n$ に対応する lepton の anomalous magnetic moment の
非標準寄与は，Ramanujan $\tau$ 関数の specific な値に比例する:
\begin{align}
\Delta a_e &\propto |\tau(1)| = 1 \quad\text{(extremely small)}\\
\Delta a_\mu &\propto |\tau(3)| = 252 \quad\text{(muon anomaly)}\\
\Delta a_\tau &\propto |\tau(5)| = 4830 \quad\text{?}
\end{align}
\end{conjecture}

もしこれが正しければ，世代構造は\textbf{Ramanujan 判別式の Fourier 展開}に
encode されている．$\Delta$ は重さ 12 cusp form で，$M_{24}$ Mathieu
群とも関連する (Mathieu Moonshine)．

これは Mathieu Moonshine を SM 世代構造に adapter する candidate path．

\subsection{貢献 5: Beilinson 予想の物理的実現}

Deligne 予想~\cite{Deligne1979}: critical motive の L 値は period と
algebraic number の積 (Deligne's period conjecture).

Wright Brothers 的形式: 物理定数は $L(s=\text{critical}, \text{motive})$
の形をとる．$\alpha^{-1}$ の leading Bernoulli 部分 $12 + 120 = 132$ は
特定の motive の $L$-value と解釈できる可能性．

\textbf{具体的挑戦}: $\alpha^{-1} - 132$ (= running correction $\approx 5.04$)
を $\zeta(5)$, $\zeta(3)$ 等の combination で書けるか? QED 高次補正の
計算結果は正に $\zeta(3), \zeta(5)$ 等を含むため，ここに rigorous な
接続が存在する可能性．

\begin{observation}[計算的目標]
$\alpha^{-1} - (2/|B_2| + 4/|B_4|) = 5.04$
を $\zeta(3), \zeta(5)$ 等の「次の次数」arithmetic invariants で
表現する．これは Beilinson 予想の物理的 test．
\end{observation}

\subsection{貢献 6: 4D--保型形式 bridge}

既存の保型形式物理は $D = 10, 11$ で主に機能．Wright Brothers は
$D = 4$ の観測量を直接 modular form に結びつける．

Examples:
\begin{itemize}
\item $\alpha^{-1}$ (4D QED) $\leftrightarrow$ $E_2, E_4$ (重さ 2, 4)
\item $\Omega_\Lambda$ (4D cosmology) $\leftrightarrow$ $E_2^2$ (重さ 4)
\item Muon $g-2$ (4D lepton) $\leftrightarrow$ $\tau(3)$ (重さ 12 cusp)
\end{itemize}

これらが正しければ，4D 物理と 2D 保型形式の\textbf{直接 bridge}が
確立．保型形式物理の scope を拡張する重要な貢献．

\subsection{貢献 7: 宇宙論定数問題の解消}

Wright Brothers の $\Omega_\Lambda = 2\pi/9$ 予測は，$10^{123}$ ズレの
問題を:
\begin{itemize}
\item $\rho_\Lambda \ll \rho_P$ の hierarchy は「UV cutoff の誤認」
\item 正しい表現は $\rho_\Lambda = c^4/(12 G \ell_H^2)$ で $\ell_H$ 依存
\item $1/12 = |B_2|/2 = |\zeta(-1)|$ が唯一の arithmetic 係数
\end{itemize}
として解消する．これは\textbf{anthropic でない，arithmetic な解決}で
landscape 問題と異なるアプローチ．

\section{新しい貢献の深掘り}

ここまでの 7 つの貢献候補のうち，最も有望な 3 つを深掘りする．

\subsection{深掘り A: 貢献 2 ($\Omega_\Lambda$ と $E_2$ quasi-modular)}

Eisenstein $E_2$ は厳密な modular form\textbf{ではない}(anomaly あり):
\begin{equation}
E_2\left(\frac{a\tau+b}{c\tau+d}\right) = (c\tau+d)^2 E_2(\tau) - \frac{6ic(c\tau+d)}{\pi}
\end{equation}
第二項の $-6ic/\pi$ が $E_2$ を modular から外している．

しかし $E_2$ の Fourier 係数は Bernoulli 数からで:
\begin{equation}
E_2(\tau) = 1 - 24\sum_{n=1}^\infty \sigma_1(n) q^n
\end{equation}
定数項 $-24 = -4/|B_2|$ (Wright Brothers には $4/|B_2|$)．

\textbf{Observation}: $E_2$ の quasi-modularity は「anomaly 的修正」で，
物理的 anomaly (量子 anomaly) と conceptually 対応する．

Wright Brothers $\Omega_\Lambda = 8\pi B_2^2$ の $B_2^2$ は:
\begin{align}
E_2^2 &= E_4 + \text{(anomaly correction)}\\
(1 - 24q + \ldots)^2 &= 1 - 48q + 576 q^2 + \ldots\\
E_4 &= 1 + 240q + \ldots
\end{align}

$E_2^2 - E_4 = -48q - 240q + O(q^2) = -288q + \ldots$?
$-288/24 = -12$. Hmm.

Precise: $E_2^2 = E_4 + (12/(\pi i))(d/d\tau) E_2$. The $d/d\tau$
makes the quasi-modular nature explicit.

\begin{conjecture}[$\Omega_\Lambda$ from $E_2$ anomaly]
宇宙論定数 $\Omega_\Lambda$ は $E_2$ の modular anomaly
(quasi-modular correction term) と関連する．具体的に，
\begin{equation}
\Omega_\Lambda \stackrel{?}{=} \frac{(\text{$E_2$ anomaly coefficient})^2}{(\text{normalization})}
\end{equation}
の形式で導かれる．
\end{conjecture}

これが正しければ，宇宙論定数は\textbf{$E_2$ の量子 anomaly の物理的
現れ}という美しい解釈になる．

\subsection{深掘り B: 貢献 4 (世代と $M_{24}$ Mathieu Moonshine)}

Mathieu moonshine~\cite{EOT2010}: K3 面の elliptic genus は
\begin{equation}
Z_{\rm K3}(\tau, z) = 24 \phi_0(\tau, z) + \text{cusp}
\end{equation}
で，係数に $M_{24}$ 既約表現次元が現れる．

Ramanujan $\tau$ 関数は $M_{24}$ とは直接関係しないが，重さ 12 cusp form
で保型形式の「類似対象」．もし lepton 世代が $M_{24}$ あるいは
Ramanujan $\tau$ と関連するなら:

\begin{itemize}
\item 3 世代 $\leftrightarrow$ $\tau(1), \tau(3), \tau(5)$ または
$M_{24}$ の 3 つの特別な conjugacy class
\item 世代混合 (PMNS) $\leftrightarrow$ $M_{24}$ 作用?
\item CP 破れ $\leftrightarrow$ modular form の複素共役?
\end{itemize}

これらは highly speculative だが，Mathieu moonshine が K3 に留まらず
SM にも関連する\textbf{extended moonshine}の可能性を示唆．
Cheng--Duncan--Harvey の Umbral Moonshine~\cite{CDH2014}は moonshine
の多様化を既に示している．

\begin{conjecture}[SM 世代 Moonshine]
Standard Model の世代構造は，Mathieu moonshine と類似の
\textbf{世代 moonshine}を持ち，Ramanujan $\tau$ 関数または関連
cusp form の Fourier coefficient が世代特有の anomaly を決定する．
\end{conjecture}

\textbf{検証方向}: $\tau$ 以外の cusp form (例: $\Delta^2$, $\Delta \cdot E_4$)
の Fourier 係数が他の SM 異常と一致するか探索．

\subsection{深掘り C: 貢献 5 (Beilinson 物理的実現)}

Deligne/Beilinson 予想は L-value が period と algebraic number の積と
予想する．具体的:
\begin{equation}
L^*(M, n) \sim (\text{period})^r \cdot (\text{algebraic number})
\end{equation}
(critical motive $M$, integer $n$)

Wright Brothers の $\alpha^{-1}$ formula:
\begin{equation}
\alpha^{-1} = \frac{2}{|B_2|} + \frac{4}{|B_4|} + g(\ldots) + d(\zeta(2d-2) - \zeta(2d-1))
\end{equation}
は $B_{2k}$ (= $\zeta(1-2k)$) と $\zeta(2k)$ の両方を含む．正と負の
integer argument で $\zeta$ を使う．これは Beilinson の「dual」構造と整合:
\begin{itemize}
\item $\zeta(1-2k)$ at negative integer: Bernoulli (arithmetic)
\item $\zeta(2k)$ at positive even: period $\pi^{2k}$
\item 両者を結ぶ: functional equation, Beilinson regulator
\end{itemize}

\begin{proposition}[Wright Brothers $\alpha^{-1}$ in Beilinson framework]
$\alpha^{-1}$ の algebraic 部分は $\zeta(-k)$, transcendental 部分
(running correction) は $\zeta(k)$ で，両者の balance が
観測値 $137.036$ を与える．
\end{proposition}

これを rigorous に書けば，Deligne--Beilinson の物理的 embodiment と
なる可能性．

\section{新しい予測}

上記の framework から派生する新しい観測可能予測:

\subsection{タウ $g-2$ の Ramanujan 予測}

\begin{conjecture}[Tau g-2 from $\tau(5)$]
タウ lepton の anomalous magnetic moment の「世代 3 寄与」は
$\tau(5) = 4830$ に比例:
\begin{equation}
\Delta a_\tau \sim 4830 \times 10^{-k}
\end{equation}
$k$ は次元解析から決まる．
\end{conjecture}

現在の tau $g-2$ 実験限界: $a_\tau \lesssim 10^{-2}$．Wright Brothers
予測 $\Delta a_\tau \sim 10^{-8}$ は検証不能．Belle II, LHCb, FCC-ee の
将来実験で検証できるか．

\subsection{Electron $g-2$ の高次補正}

$a_e$ の 5-loop 計算には $\zeta(3), \zeta(5), \zeta(7)$ が登場．
Wright Brothers は:
\begin{equation}
a_e^{(2n)} \sim (\text{QED factor}) \times (\text{arithmetic combination of } \zeta(2k+1))
\end{equation}
を予測する．これは既存の QED 計算と matching するはず．
\textbf{既存計算結果}との詳細比較は future work．

\subsection{Higgs self-coupling $\lambda_H$}

$\lambda_H \approx 0.129$ は現在実験から決定されている．Wright Brothers
arithmetic 表現の candidate:
\begin{itemize}
\item $\lambda_H \stackrel{?}{=} B_2/\pi = (1/6)/\pi \approx 0.053$: 違う
\item $\lambda_H \stackrel{?}{=} \pi B_2 = \pi/6 \approx 0.524$: 違う
\item $\lambda_H \stackrel{?}{=} 1/(2\pi^2) \approx 0.0507$: 近い桁
\end{itemize}

Clean な表現は現時点で見つからない．Higgs は EWSB に関連する
dynamical パラメータで，zero-point arithmetic でなく RG evolution
由来かもしれない．

\subsection{宇宙論 $H_0$ tension と $\Omega_\Lambda$}

Wright Brothers 予測 $\Omega_\Lambda = 2\pi/9 = 0.6981$ は
$H_0$ に独立．Planck 2018 ($H_0 = 67.66$) で $\Omega_\Lambda = 0.685$,
SH0ES ($H_0 = 73.0$) で異なる値．

もし $\Omega_\Lambda = 2\pi/9$ が正確なら，$H_0$ tension は:
\begin{itemize}
\item Planck 側が若干過小評価 ($0.685 \to 0.698$)
\item または SH0ES 側が正しく，他のパラメータが shift
\end{itemize}

2027--2028 の次世代観測で決着．

\section{Wright Brothers の「立ち位置」}

\subsection{既存分野との住み分け}

Wright Brothers がoccupy する niche:

\begin{center}
\begin{tabular}{lll}
\toprule
既存分野 & 扱う物理 & WB の相補性 \\
\midrule
Heterotic string & 10D partition & WB: 4D constants \\
Monstrous Moonshine & Monster VOA & WB: SM observables \\
Mathieu Moonshine & K3 geometry & WB: lepton generations (conj.)\\
GV/DT invariants & BPS counting & WB: continuous coupling \\
BH microstates & Igusa forms & WB: vacuum energy \\
Connes--Marcolli & Spectral triples & WB: Arithmetic values \\
\bottomrule
\end{tabular}
\end{center}

Wright Brothers は「counting theory」ではなく「constant determination
theory」として位置付けられる．既存分野と\textbf{overlap せず補完}する．

\subsection{方法論的相違}

\begin{itemize}
\item \textbf{既存分野}: 数学的 consistency (modular invariance, $SL_2$ 共変)
から出発 → 物理的対象 (partition function) を特定
\item \textbf{Wright Brothers}: 観測された数値 (コミュ値) から出発 →
arithmetic 表現を探索 → 数学的 framework に embed
\end{itemize}

順序が逆．これは\textbf{Einstein 的 vs Dirac 的}と言える:
\begin{itemize}
\item Einstein: 原理から演繹 (一般共変性 → GR)
\item Dirac: 美から演繹 (対称性 → Dirac eqn)
\item Wright Brothers: 観測的数値一致 → 原理探索
\end{itemize}

Wright Brothers はより Dirac 的で，その方法論にはリスクがある
(LNH 的失敗の可能性)．

\subsection{検証の roadmap}

今後 5 年で Wright Brothers が\textbf{Einstein 側に移動する}ための
条件:

\begin{enumerate}
\item \textbf{Pre-diction の文書化}: 未観測量 (tau $g-2$, dark matter
mass 等) の arithmetic 予測を事前登録
\item \textbf{Independent experimental verification}:
\begin{itemize}
\item Euclid, DESI, LSST ($\Omega_\Lambda$ 0.5\%)
\item Belle II, LHCb (tau $g-2$)
\item 水晶 BAW 実験 ($a_c \approx 88\,\mu$m)
\end{itemize}
\item \textbf{理論的導出の rigorous 化}: CKN bound 係数の $1/12$,
Wheeler--DeWitt DN 境界の厳密定式化, Kreimer Hopf 代数 + Wright
Brothers 式の統合
\item \textbf{既存分野との bridge}: Heterotic string, Moonshine と
の explicit な接続計算
\item \textbf{失敗の honest report}: 合わない予測も公開し selection bias
を minimize
\end{enumerate}

これらが実現すれば，Wright Brothers は既存の保型形式物理に
\textbf{実質的な補完}を提供する research program となる．

\section{結論}

本稿の主要 claim:

\begin{enumerate}
\item 既存の保型形式物理 (heterotic, moonshine, topological strings, BH)
は深い数学的成果を多数持つが，\textbf{4D 低エネルギー SM パラメータの
決定}という領域では何もできていない
\item この空白は Wright Brothers プロジェクトの natural niche:
$\alpha^{-1} = 132 + $ running, $\Omega_\Lambda = 2\pi/9$, muon $g-2 = 252$
等の arithmetic 表現は既存 program には欠けている
\item Wright Brothers は保型形式物理の\textbf{補完}として，
「counting」ではなく「determination」を担う
\item 具体的な bridge: $E_2$ quasi-modular と $\Omega_\Lambda$,
Ramanujan $\tau$ と世代構造, Beilinson 予想と $\alpha^{-1}$ loop 展開
\item Einstein 側に移動するには pre-diction 文書化，実験検証，
rigorous 理論化が必要．2028 年頃までに決着する可能性
\end{enumerate}

これが実現すれば，物理定数の起源を Spec$(\mathbb{Z})$ arithmetic から
第一原理的に導出する新しい programm が確立する．既存の保型形式物理は
Wright Brothers の「重要な上位構造」として，SM と GR の統合に
貢献する．

\begin{thebibliography}{99}

\bibitem{GSW1987}
M.~B.~Green, J.~H.~Schwarz, E.~Witten, \textit{Superstring Theory}
(Cambridge Univ.\ Press, 1987).

\bibitem{Borcherds1992}
R.~Borcherds, Invent.\ Math.\ \textbf{109}, 405 (1992).

\bibitem{FLM1988}
I.~Frenkel, J.~Lepowsky, A.~Meurman, \textit{Vertex Operator
Algebras and the Monster} (Academic Press, 1988).

\bibitem{EOT2010}
T.~Eguchi, H.~Ooguri, Y.~Tachikawa,
Exp.\ Math.\ \textbf{20}, 91 (2011).

\bibitem{Gannon2016}
T.~Gannon, Adv.\ Math.\ \textbf{301}, 322 (2016).

\bibitem{GV1998}
R.~Gopakumar, C.~Vafa, hep-th/9809187, hep-th/9812127 (1998).

\bibitem{VafaWitten1994}
C.~Vafa, E.~Witten, Nucl.\ Phys.\ B \textbf{431}, 3 (1994).

\bibitem{DMVV1997}
R.~Dijkgraaf, G.~Moore, E.~Verlinde, H.~Verlinde,
Commun.\ Math.\ Phys.\ \textbf{185}, 197 (1997).

\bibitem{Dabholkar2012}
A.~Dabholkar, S.~Murthy, D.~Zagier, arXiv:1208.4074.

\bibitem{BCOV1994}
M.~Bershadsky, S.~Cecotti, H.~Ooguri, C.~Vafa,
Commun.\ Math.\ Phys.\ \textbf{165}, 311 (1994).

\bibitem{Susskind2003}
L.~Susskind, hep-th/0302219.

\bibitem{Weinberg1987}
S.~Weinberg, Phys.\ Rev.\ Lett.\ \textbf{59}, 2607 (1987).

\bibitem{KKLT2003}
S.~Kachru, R.~Kallosh, A.~Linde, S.~Trivedi,
Phys.\ Rev.\ D \textbf{68}, 046005 (2003).

\bibitem{Vafa2005}
C.~Vafa, hep-th/0509212.

\bibitem{Deligne1979}
P.~Deligne, Proc.\ Symp.\ Pure Math.\ \textbf{33.2}, 313 (1979).

\bibitem{Beilinson1984}
A.~Beilinson, J.\ Soviet Math.\ \textbf{30}, 2036 (1985).

\bibitem{BlochKato1990}
S.~Bloch, K.~Kato, \textit{The Grothendieck Festschrift}, Vol.\ I
(Birkhäuser, 1990).

\bibitem{Kreimer1999}
D.~Kreimer, \textit{Knots and Feynman Diagrams}
(Cambridge Univ.\ Press, 2000).

\bibitem{CDH2014}
M.~Cheng, J.~Duncan, J.~Harvey,
Commun.\ Number Theory Phys.\ \textbf{8}, 101 (2014).

\bibitem{wb-spec-z}
Wright Brothers Research Program, ``Spec$(\mathbb{Z})$ alphabet 深掘り'',
\texttt{spec\_z\_alphabet\_deeper\_ja.tex} (2026).

\bibitem{wb-alpha}
Wright Brothers Research Program, ``微細構造定数の算術的表現'',
\texttt{alpha\_arithmetic\_ja.tex} (2026).

\bibitem{wb-wdw}
Wright Brothers Research Program, ``WDW 時間的 Dirichlet--Neumann
真空と $\Omega_\Lambda = 2\pi/9$'',
\texttt{wdw\_temporal\_dn\_ja.tex} (2026).

\end{thebibliography}

\end{document}
